\documentclass[11pt]{article}
\usepackage{series}

\title{Finite Coherent Overlap Vertices and Their Local Contact Limit\\
\large Conditional EFT Interpretation and the Renormalization Boundary}
\author{Peter Nero}
\date{July 2026, Version 1}

\begin{document}
\maketitle

\begin{abstract}
Local quantum field theory writes interactions as pointwise products, for example
\(\int_X \lambda \phi(x)^4\,dx\).  In perturbation theory these contact vertices generate
coincidence limits, momentum-conservation delta functions, and ultraviolet singularities.
Building on the preceding papers on Dirac deltas, coherent Green functions, gauge fixing, and
measurement effects, this paper develops the contact-interaction member of the same
projection-kernel program.  The central claim is narrow: a point-local interaction vertex is the
zero-width limit of a finite coherent overlap vertex.

The mathematical core is an approximate-identity theorem.  Given a bounded smoothing
kernel \(K_\epsilon(x,y)\), define the smeared field
\(\phi_\epsilon(x)=\int_X K_\epsilon(x,y)\phi(y)\,dy\) and the finite overlap interaction
\[
S_{\mathrm{int},\epsilon}[\phi]
  = \frac{\lambda}{4!}\int_X \phi_\epsilon(x)^4\,dx .
\]
For smooth fields, \(S_{\mathrm{int},\epsilon}\to S_{\mathrm{int}}=
\frac{\lambda}{4!}\int_X\phi(x)^4\,dx\) as \(\epsilon\downarrow0\).  On a compact spectral
model, replacing the identity by a finite spectral projector \(\Pi_\Lambda\) yields a finite
vertex tensor and finite coincident covariance in the retained sector.  In momentum space,
heat-kernel smearing multiplies propagators and vertices by rapidly decaying form factors,
turning point-contact singularities into finite coherent-overlap expressions before the sharp
limit is taken.

The MTT interpretation is conditional.  If a selected physical sector supplies such a finite
overlap kernel and an intertwiner to external fields, the local vertex may be read as its sharp
limit.  At nonzero width the construction is generally a nonlocal effective interaction, not
automatically a local ultraviolet completion.  A local parent mediator or an equivalent
localization theorem would be needed to change that classification.  The paper proves
kernel convergence, the general \(n\)-point sharp limit, and finite-cutoff statements.  It
does not infer all-loop renormalizability, ultraviolet completion, Lorentzian causality,
unitarity, or gauge/BRST consistency from finite width alone.
\end{abstract}

\section*{Version 1 Revision Note}
\begin{description}
\item[Supersedes] The unversioned April 2026 manuscript with the same subject.
\item[Reason] The earlier text correctly proved an approximate-identity limit but could be
read as promoting finite overlap directly to a physical MTT interaction.
\item[Resolution] This version classifies the finite-width vertex as a nonlocal EFT
interaction unless a local parent or localization theorem is supplied, and separates the
proved finite-cutoff statements from all-loop, causal, unitary, and gauge claims.
\item[Retained result] The sharp-limit theorem and finite spectral vertex tensors are
unchanged.
\item[Remaining boundary] A selected external-field source map and physical compatibility
certificates are still required.
\end{description}

\section{Purpose and claim discipline}

The preceding papers established four projection-kernel replacements:
\[
\delta(x-y)\leadsto K_{\mathrm{coh}}(x,y),
\qquad
LG=\delta\leadsto LG_{\mathrm{coh}}=K_{\mathrm{coh}},
\]
\[
\delta(G[A])\leadsto \mathcal K_{\epsilon}(G[A]),
\qquad
|x\rangle\langle x|\leadsto E_x^\epsilon .
\]
The present paper treats the next major occurrence of delta-like structure: contact
interactions.

In ordinary local field theory, an interaction such as
\[
S_{\mathrm{int}}[\phi]=\frac{\lambda}{4!}\int_X \phi(x)^4\,dx
\]
is point-local.  In perturbation theory this point-locality appears as coincidence limits,
delta functions enforcing exact vertex bookkeeping, and ultraviolet divergences.

The central thesis is:
\[
\boxed{\text{A contact interaction is the zero-width limit of a finite coherent overlap vertex.}}
\]

\subsection{Non-claims}

This paper does not claim that renormalization is unnecessary, that all quantum field theories
become finite by the same smearing rule, or that a complete Standard Model calculation is
performed here.  The claims are narrower.

\begin{center}
\begin{tabular}{p{0.31\linewidth}p{0.58\linewidth}}
\hline
\textbf{Level} & \textbf{Claim}\\
\hline
Proved & Smeared overlap vertices converge to local contact vertices on smooth fields.\\
Proved & Finite spectral projection yields finite-dimensional vertex tensors and finite
coincident covariance in the retained sector.\\
Standard analytic fact & Heat-kernel or spectral damping improves ultraviolet behavior in loop
integrals.\\
MTT interpretation & Local contact vertices are singular downstream shadows of finite coherent
overlap.\\
Research program & Renormalization can be partly re-read as repair of over-sharp projection.\\
\hline
\end{tabular}
\end{center}

\subsection{Physical admissibility caveat}

A finite-width vertex is not automatically a physically acceptable replacement for a local
interaction.  Generic nonlocal smearing can spoil microcausality, reflection positivity,
unitarity, Ward identities, or BRST consistency.  The present paper therefore distinguishes
two notions:
\[
\text{arbitrary regulator kernel}
\qquad\text{versus}\qquad
\text{admissible coherent projection kernel}.
\]
Only the second has the intended MTT status.  An admissible kernel must be generated by the
same coherent-sector projection, overlap geometry, and symmetry constraints that define the
effective theory.  In gauge theories, for example, finite-width vertices must preserve the
relevant Ward or Slavnov--Taylor identities.  In Lorentzian settings, they must also be
compatible with causal propagation on the admissible slab.

Thus the paper proves convergence and finiteness statements for controlled kernels, but does
not claim that arbitrary smearing produces a consistent QFT.

\begin{remark}[EFT classification]
At fixed nonzero width, \(V_\epsilon\) couples fields at separated points.  It is therefore a
nonlocal EFT vertex unless it is obtained by integrating out a declared local mediator or a
localization theorem supplies an equivalent local parent theory.  Finiteness of a vertex
tensor at fixed cutoff does not imply all-loop renormalizability or ultraviolet completion.
\end{remark}

\section{Local contact vertices and their hidden delta structure}

A local quartic vertex can be written formally as
\[
S_{\mathrm{int}}[\phi]
=\frac{\lambda}{4!}\int_X \phi(x)^4\,dx .
\]
Equivalently, it is the four-field functional
\[
S_{\mathrm{int}}[\phi]
=
\frac{\lambda}{4!}
\int_{X^4}
V_{\delta}(x_1,x_2,x_3,x_4)
\phi(x_1)\phi(x_2)\phi(x_3)\phi(x_4)
\,dx_1\cdots dx_4,
\]
where the contact vertex distribution is
\[
V_{\delta}(x_1,x_2,x_3,x_4)
=
\int_X
\prod_{j=1}^{4}\delta(x-x_j)\,dx .
\]
Thus the local vertex already contains a product of selection kernels.  It says that all four
fields interact at exactly the same point.  The MTT diagnostic asks: what finite coherent
overlap has been idealized as this exact coincidence?

\section{Finite coherent overlap vertices}

Let \(K_\epsilon(x,y)\) be a family of smooth kernels on a compact Riemannian manifold
\((X,g)\).  Assume \(K_\epsilon\) is an approximate identity: for every smooth \(f\),
\[
K_\epsilon f(x):=\int_X K_\epsilon(x,y)f(y)\,dy\longrightarrow f(x)
\]
in \(C^\infty(X)\) as \(\epsilon\downarrow0\).  Examples include heat kernels
\(H_\epsilon=e^{-\epsilon\Delta}\) and spectral cutoffs approaching the identity.

Define the smeared field
\[
\phi_\epsilon(x)=K_\epsilon\phi(x)
=\int_XK_\epsilon(x,y)\phi(y)\,dy .
\]
The finite coherent overlap vertex is
\[
S_{\mathrm{int},\epsilon}[\phi]
=
\frac{\lambda}{4!}\int_X \phi_\epsilon(x)^4\,dx .
\]
Equivalently,
\[
S_{\mathrm{int},\epsilon}[\phi]
=
\frac{\lambda}{4!}
\int_{X^4}
V_\epsilon(x_1,x_2,x_3,x_4)
\prod_{j=1}^4\phi(x_j)\,
dx_1\cdots dx_4,
\]
where
\[
V_\epsilon(x_1,x_2,x_3,x_4)
=
\int_X\prod_{j=1}^{4}K_\epsilon(x,x_j)\,dx .
\]
For finite \(\epsilon\), the vertex no longer forces exact coincidence.  It couples fields
whose support overlaps within the coherent resolution scale of \(K_\epsilon\).

\section{Theorem: overlap vertices converge to contact vertices}

\begin{theorem}[Contact vertex as sharp-overlap limit]
Let \(X\) be compact and let \(K_\epsilon\) be an approximate identity on \(C^\infty(X)\).
For every \(\phi\in C^\infty(X)\),
\[
S_{\mathrm{int},\epsilon}[\phi]\longrightarrow
S_{\mathrm{int}}[\phi]
=
\frac{\lambda}{4!}\int_X\phi(x)^4\,dx
\]
as \(\epsilon\downarrow0\).
\end{theorem}

\begin{proof}
Since \(K_\epsilon\phi\to\phi\) in \(C^\infty(X)\), in particular
\(K_\epsilon\phi\to\phi\) uniformly.  Write \(\phi_\epsilon=K_\epsilon\phi\).  Then
\[
|\phi_\epsilon^4-\phi^4|
=
|\phi_\epsilon-\phi|\,
|\phi_\epsilon^3+\phi_\epsilon^2\phi+\phi_\epsilon\phi^2+\phi^3|.
\]
Uniform convergence implies \(\sup_X|\phi_\epsilon-\phi|\to0\), while
\(\phi_\epsilon\) and \(\phi\) are uniformly bounded for sufficiently small \(\epsilon\).
Hence \(\phi_\epsilon^4\to\phi^4\) uniformly.  Since \(X\) has finite volume,
\[
\int_X \phi_\epsilon(x)^4\,dx \to \int_X\phi(x)^4\,dx .
\]
Multiplying by \(\lambda/4!\) gives the result.
\end{proof}

\begin{corollary}[Contact distribution as vertex limit]
In the distributional sense on \(X^4\),
\[
V_\epsilon(x_1,x_2,x_3,x_4)\to
V_\delta(x_1,x_2,x_3,x_4)
=
\int_X\prod_{j=1}^4\delta(x-x_j)\,dx .
\]
\end{corollary}

\begin{proof}
For a test function of product form
\(F(x_1,x_2,x_3,x_4)=\prod_{j=1}^4\phi_j(x_j)\), the pairing with \(V_\epsilon\) is
\[
\int_X\prod_{j=1}^4(K_\epsilon\phi_j)(x)\,dx,
\]
which converges to
\[
\int_X\prod_{j=1}^4\phi_j(x)\,dx,
\]
the pairing with \(V_\delta\).  Finite linear combinations of product test functions are
dense in the usual test-function topology on compact products; continuity gives the general
case.
\end{proof}

\section{General \texorpdfstring{\(n\)}{n}-point overlap vertices}

The quartic case is not special.  For an \(n\)-field local vertex, define
\[
V_{\delta}^{(n)}(x_1,\ldots,x_n)
=
\int_X\prod_{j=1}^{n}\delta(x-x_j)\,dx ,
\]
so that the corresponding local interaction is
\[
S_{\mathrm{int}}^{(n)}[\phi]
=
\frac{\lambda_n}{n!}
\int_{X^n}
V_{\delta}^{(n)}(x_1,\ldots,x_n)
\prod_{j=1}^{n}\phi(x_j)\,dx_1\cdots dx_n
=
\frac{\lambda_n}{n!}\int_X\phi(x)^n\,dx .
\]
The finite overlap replacement is
\[
V_{\epsilon}^{(n)}(x_1,\ldots,x_n)
=
\int_X\prod_{j=1}^{n}K_\epsilon(x,x_j)\,dx ,
\]
and
\[
S_{\mathrm{int},\epsilon}^{(n)}[\phi]
=
\frac{\lambda_n}{n!}\int_X (K_\epsilon\phi)(x)^n\,dx .
\]

\begin{proposition}[General \(n\)-point sharp-overlap limit]
Let \(X\) be compact, let \(K_\epsilon\) be an approximate identity on \(C^\infty(X)\), and
let \(n\ge2\).  For every \(\phi\in C^\infty(X)\),
\[
S_{\mathrm{int},\epsilon}^{(n)}[\phi]
\longrightarrow
S_{\mathrm{int}}^{(n)}[\phi]
=
\frac{\lambda_n}{n!}\int_X\phi(x)^n\,dx .
\]
Moreover \(V_{\epsilon}^{(n)}\to V_{\delta}^{(n)}\) in distributions on \(X^n\).
\end{proposition}

\begin{proof}
The functional convergence follows from \(K_\epsilon\phi\to\phi\) uniformly and the identity
\[
a^n-b^n=(a-b)\sum_{r=0}^{n-1}a^{n-1-r}b^r .
\]
Uniform boundedness of \(K_\epsilon\phi\) for small \(\epsilon\) gives uniform convergence of
\((K_\epsilon\phi)^n\) to \(\phi^n\), and compactness of \(X\) permits integration.

For the distributional statement, pair \(V_{\epsilon}^{(n)}\) with product test functions
\(\prod_{j=1}^n\phi_j(x_j)\).  The pairing is
\[
\int_X\prod_{j=1}^n(K_\epsilon\phi_j)(x)\,dx,
\]
which converges to
\[
\int_X\prod_{j=1}^n\phi_j(x)\,dx.
\]
Density of finite sums of product test functions in the test-function topology on \(X^n\)
extends the result.
\end{proof}

This general form is the useful one for perturbation theory: every point-local \(n\)-leg
vertex is the sharp limit of a finite coherent \(n\)-fold overlap.

\section{Finite spectral projection and finite vertex tensors}

A particularly clean model uses a finite spectral projector.  Let
\[
\Pi_\Lambda f
=
\sum_{\lambda_n\le\Lambda}\langle\phi_n,f\rangle\phi_n
\]
be the spectral projector for a nonnegative Laplace-type operator on compact \(X\).  Define
\[
\phi_\Lambda=\Pi_\Lambda\phi .
\]
The projected quartic interaction is
\[
S_{\mathrm{int},\Lambda}[\phi]
=
\frac{\lambda}{4!}\int_X(\Pi_\Lambda\phi)(x)^4\,dx .
\]
Expanding
\[
\Pi_\Lambda\phi=\sum_{\lambda_a\le\Lambda}\phi_a\,e_a(x)
\]
gives
\[
S_{\mathrm{int},\Lambda}
=
\frac{\lambda}{4!}
\sum_{a,b,c,d\in I_\Lambda}
V_{abcd}\,\phi_a\phi_b\phi_c\phi_d,
\]
where \(I_\Lambda=\{a:\lambda_a\le\Lambda\}\) and
\[
V_{abcd}=\int_X e_a(x)e_b(x)e_c(x)e_d(x)\,dx .
\]
For finite \(\Lambda\), this is a finite-dimensional tensor.  There are no coincident-point
products of distributions inside the retained sector.

\begin{proposition}[Finite coincident covariance in a finite coherent sector]
Let \(C_\Lambda=\Pi_\Lambda L^{-1}\Pi_\Lambda\), where \(L\) is strictly positive on
\(L^2(X)\).  Then \(C_\Lambda\) has smooth finite-rank kernel
\[
C_\Lambda(x,y)=\sum_{\lambda_n\le\Lambda}\frac{e_n(x)e_n^\ast(y)}{\mu_n},
\]
where \(L e_n=\mu_n e_n\) in the diagonal case.  In particular, \(C_\Lambda(x,x)<\infty\)
for every \(x\in X\).
\end{proposition}

\begin{proof}
The sum contains finitely many smooth terms, each with finite coefficient \(1/\mu_n\).
Therefore \(C_\Lambda(x,y)\) is smooth and finite-rank.  Evaluating at \(x=y\) gives a finite
sum of finite numbers.
\end{proof}

This proposition is the rigorous finite-cutoff version of the slogan:
\[
\boxed{\text{coincident-point singularities arise only after the coherent cutoff is removed}.}
\]

\section{Momentum-space model on the flat torus}

Let \(X=\mathbb T^d\) and use Fourier modes
\[
e_k(x)=\frac{1}{(2\pi)^{d/2}}e^{ik\cdot x},
\qquad k\in\mathbb Z^d.
\]
For \(L_m=-\Delta+m^2\), the free covariance is
\[
C(k)=\frac{1}{|k|^2+m^2}.
\]

A sharp spectral cutoff retains \(|k|\le\Lambda\).  The projected covariance is
\[
C_\Lambda(x,y)=\frac{1}{(2\pi)^d}
\sum_{|k|\le\Lambda}
\frac{e^{ik\cdot(x-y)}}{|k|^2+m^2}.
\]
For finite \(\Lambda\), \(C_\Lambda(x,x)\) is finite:
\[
C_\Lambda(x,x)=\frac{1}{(2\pi)^d}
\sum_{|k|\le\Lambda}
\frac{1}{|k|^2+m^2}<\infty .
\]
As \(\Lambda\to\infty\), this diagonal quantity diverges for \(d\ge2\), logarithmically in
\(d=2\) and by power law in \(d\ge3\).  This is the standard ultraviolet coincidence problem.

A heat-filtered version replaces the sharp cutoff by
\[
C_\tau(k)=\frac{e^{-\tau |k|^2}}{|k|^2+m^2}.
\]
Then
\[
C_\tau(x,x)=\frac{1}{(2\pi)^d}
\sum_{k\in\mathbb Z^d}
\frac{e^{-\tau |k|^2}}{|k|^2+m^2}<\infty
\]
for every \(\tau>0\).  The divergence reappears only as \(\tau\downarrow0\).

Thus the finite coherent kernel does not ``solve'' the local theory by magic.  It displays
exactly where the singularity enters: at the zero-width limit.

\section{Loop integrals and over-sharp projection}

In Euclidean momentum space, the one-loop tadpole in a local scalar theory contains the
formal integral
\[
I_{\mathrm{local}}=\int_{\mathbb R^d}\frac{d^dk}{(2\pi)^d}\frac{1}{k^2+m^2},
\]
which diverges for \(d\ge2\).  A heat-coherent replacement gives
\[
I_\tau=\int_{\mathbb R^d}\frac{d^dk}{(2\pi)^d}
\frac{e^{-\tau k^2}}{k^2+m^2}.
\]
For every \(\tau>0\), \(I_\tau<\infty\).  The local divergent integral is recovered as
\(\tau\downarrow0\).

This illustrates the structural point:
\[
\boxed{\text{UV divergence is often the cost of taking a coherent overlap width to zero}.}
\]
The statement is intentionally limited.  It does not replace the full renormalization program.
It explains why renormalization is encountered precisely where local field theory forms
products or loops at exact coincidence.

\section{Regulator versus coherent projection kernel}

Mathematically, the kernels introduced above resemble familiar regulators.  A heat factor
\(e^{-\tau k^2}\), a spectral cutoff \(\Pi_\Lambda\), or a smooth compactly supported kernel can
all improve ultraviolet behavior.  The MTT interpretation is not that regulators are new.
Rather, it changes the status of the finite object.

\[
\begin{array}{c|c}
\text{Ordinary regulator use} & \text{MTT coherent-kernel reading}\\
\hline
\text{auxiliary computational device} & \text{finite admissible overlap structure}\\
\text{removed after renormalization} & \text{may encode physical coherence scale}\\
\text{not usually part of ontology} & \text{downstream shadow of projection}\\
\text{chosen for convenience} & \text{constrained by coherent-sector geometry}
\end{array}
\]

Therefore the claim is not simply ``add a cutoff.''  The claim is that in an MTT-compatible
description the finite kernel is the prior object, and the local contact vertex is the idealized
limit.  Ordinary renormalization remains necessary when one insists on taking that local
limit.

\section{Renormalization as repair of over-sharp projection}

From the present viewpoint, renormalization has a two-sided interpretation.

First, it is the standard and indispensable method for defining local QFT after singular
limits have been taken.  Nothing in this paper challenges the causal-local or renormalization-
group constructions \cite{EpsteinGlaser1973,WilsonKogut1974}.

Second, it can be structurally re-read as a repair mechanism:
\[
\begin{array}{c}
\text{finite overlap vertex}\longrightarrow\text{zero-width contact vertex}\\
\phantom{\text{finite overlap vertex}}\longrightarrow
\text{coincident singularities}\longrightarrow\text{renormalization}
\end{array}
\]
In MTT terms, the local theory has projected away the finite overlap structure and then must
restore predictivity by absorbing the consequences into parameters, counterterms, or running
couplings.

Thus:
\begin{center}
\fbox{\parbox{0.86\linewidth}{\centering
Renormalization repairs the downstream theory after admissible overlap has been
idealized as point contact.}}
\end{center}
This is not a new computational scheme by itself.  It is a structural reinterpretation of why
the scheme is needed.  In particular, the existence of a finite kernel does not abolish
renormalization in the sharp local limit.  It identifies the finite object whose removal creates
the singular problem that renormalization then controls.

\section{MTT interpretation}

The MTT reading is direct.

\begin{center}
\begin{tabular}{p{0.34\linewidth}p{0.55\linewidth}}
\hline
\textbf{Standard QFT object} & \textbf{MTT reading}\\
\hline
\(\delta(x-y)\) & singular identity kernel\\
\(\phi(x)^4\) & zero-width contact overlap\\
\(V_\delta\) & exact coincidence selection among four fields\\
\(V_\epsilon\) & finite coherent overlap vertex\\
finite \(\Lambda\) or \(\tau>0\) & retained coherent sector / admissible resolution\\
coincident singularity & failure of zero-width contact idealization\\
counterterms/running & repair of over-sharp projection in the effective theory\\
\hline
\end{tabular}
\end{center}

The native MTT object is not the local vertex \(V_\delta\), but the finite overlap vertex
\(V_{\mathrm{coh}}\).  The local vertex is the singular downstream limit.

\section{Scope of the physical claim}

The mathematical statements above are deliberately weaker than the possible physical
interpretation.  What is proved is that contact vertices arise as sharp limits of finite overlap
vertices and that finite spectral sectors remove coincident distributional products inside the
retained sector.  What is not proved here is that a particular experimentally correct quantum
field theory is obtained by a chosen kernel.

The physical MTT claim is conditional:

\begin{center}
\fbox{\parbox{0.86\linewidth}{\centering
If an effective local QFT is the sharp limit of an admissible coherent projection, then its
contact vertices should be read as zero-width overlap limits.}}
\end{center}

This conditional form is important.  It avoids the false inference that any smoothing kernel is
physically meaningful.  Only kernels tied to the admissibility, symmetry, and projection
structure of the underlying theory qualify as coherent kernels in the MTT sense.

\section{Relation to the previous papers}

The present paper extends the sequence as follows.

\begin{center}
\begin{tabular}{p{0.34\linewidth}p{0.55\linewidth}}
\hline
\textbf{Previous paper} & \textbf{Replacement}\\
\hline
Delta foundation & \(\delta(x-y)\leadsto K_{\mathrm{coh}}(x,y)\)\\
Green functions & \(LG=\delta\leadsto LG_{\mathrm{coh}}=K_{\mathrm{coh}}\)\\
Gauge fixing & \(\delta(G[A])\leadsto\) finite gauge tube\\
Measurement & \(|x\rangle\langle x|\leadsto\) finite POVM effect\\
This paper & \(\phi(x)^4\leadsto\) finite coherent overlap vertex\\
\hline
\end{tabular}
\end{center}

The common pattern is now visible:
\begin{center}
\fbox{\parbox{0.78\linewidth}{\centering
singular local object \(=\) zero-width limit of a bounded admissible kernel.}}
\end{center}

\section{Scope of proof}

The proof content of this paper is deliberately modest.

\begin{center}
\begin{tabular}{p{0.38\linewidth}p{0.50\linewidth}}
\hline
\textbf{Claim} & \textbf{Status}\\
\hline
Smeared smooth vertices converge to contact vertices & proved\\
Finite spectral sectors have finite vertex tensors & proved\\
Finite spectral covariance has finite diagonal & proved\\
Heat damping makes tadpole-type integrals finite & standard analytic fact, illustrated\\
General \(n\)-point contact vertices arise as sharp overlap limits & proved for smooth fields\\
Arbitrary smearing preserves locality, unitarity, and gauge symmetry & not claimed\\
All QFT renormalization is solved by MTT kernels & not claimed\\
Finite contact width proves all-loop UV completion & not claimed\\
The Standard Model has been recalculated with these kernels & not claimed\\
\hline
\end{tabular}
\end{center}

\section{Conclusion}

The local quartic vertex is the zero-width limit of a finite overlap vertex.  At finite
spectral bandwidth the vertex tensor and coincident covariance are finite in the retained
sector.

The MTT reading is conditional: fixed width defines a nonlocal EFT until locality, causal
support, unitarity, and gauge/BRST compatibility are proved.  Renormalization remains the
standard sharp-limit analysis.

\begin{thebibliography}{9}

\bibitem{EpsteinGlaser1973}
H. Epstein and V. Glaser,
``The role of locality in perturbation theory,''
\emph{Annales de l'Institut Henri Poincar\'e A} \textbf{19} (1973), 211--295.
\url{https://eudml.org/doc/75788}

\bibitem{WilsonKogut1974}
K. G. Wilson and J. Kogut,
``The renormalization group and the \(\epsilon\) expansion,''
\emph{Physics Reports} \textbf{12} (1974), 75--199.
\url{https://doi.org/10.1016/0370-1573(74)90023-4}

\end{thebibliography}

\end{document}
